\section{Analytical realization still required}
The change from $Y_R$ to $Y_R^*$ is explicit, and every periodic record used
by Theorem~\ref{thm:joint-periodic} has been recomputed for the actual first
return. The complete periodic annihilator is now trivial, uniformly in the
sense of Corollary~\ref{cor:compact-periodic}. The former nonzero-roof
periodic obstruction is therefore resolved for this section. Moreover,
Theorem~\ref{thm:log-period-separation} supplies an explicit large-frequency
period discrepancy and Corollary~\ref{cor:approximate-phase-bound} excludes
arbitrarily accurate continuous unit-modulus approximate phases at the
printed rate. This does not identify the section with a Young magnet,
construct its weighted operator space, or regularize a merely measurable
spectral phase.

Three analytical tasks remain before a full raw LLT is released. First,
construct a common strong/weak realization of the actual unbounded induced
observables, with parameter control and enough regularity to evaluate
spectral coboundaries at the regular periodic points. Second, establish
the variance--coboundary criterion, covariance continuity and the returned
temporal estimates required for large roof frequencies. The real period
rank and compact periodic separation are the inputs for these steps, not
substitutes for their functional analysis. Third, sum the central critical
and singular branch contributions and prove an integrable residual with
\eqref{eq:local-tail}, or the stronger criterion \eqref{eq:tailcriterion}.
Remote edge mass may be treated by Corollary~\ref{cor:remote-edges}; this
does not settle the central branch sum.

The inequalities in Lemma~\ref{lem:splice} must be checked with actual
derivative-growth constants. The weighted versions for
Proposition~\ref{prop:conditioning} require their own uniform insertion
class. Exact induced means, fixed-record density and first-moment continuity,
and uniform stationary finite-endpoint control are now proved. The
Gaussian-scale process clock statement and uniform parameter Green--Kubo
comparison retain their printed hypotheses. The first-order functional
clock and the sublinear maximum over a growing observation interval are
now supplied by Theorem~\ref{thm:uniform-physical-clock}; neither is
substituted for those Gaussian-scale hypotheses. No full A3
conditioning input is exported until these requirements are established.

The original raw mixed-density local limit and its geometric error
propagation remain the target. The new period family, joint phase theorem,
local inversion identity and exact mean formulas are available independently
of the unresolved global estimates. The old section's nonintegrable raw
transform is not reused as an integrable transform for the new section.

\subsection*{Source and verification}
The A2 geometric submission is frozen at commit
\begin{center}\small\texttt{4557df22f5c72bc80943690ecd6c2e39de3734ab}.\end{center} The first A2-DYN
source is v1 at
\begin{center}\small\texttt{36f1365041de95ca478739a1e2984734c72f95aa}.\end{center}
The immediate revision base is v3 at
\begin{center}\small\texttt{55ab80ed1a4f11b58ce9a88c365b0cea2e2ce191}.\end{center}
All 28 proof bodies and 78 labels of that source remain active. Nine
proved results are added in the quantitative-period and uniform-clock
sections. The locally delivered v2 manuscript is also retained in full at the level
of all proof bodies and labels. The v1 proof bodies and labels are retained. The rational winding certificate is unchanged.
The new minimum certificate uses rational interval evaluation and a global
strong-convexity inequality at finitely many parameters and periods. The
infinite period-excess theorem is proved analytically, not by extrapolating
those samples. Independent finite mechanical checks import no author solver.
AI-assisted analysis contributed to this manuscript. An independent human
specialist review and a formal continuum proof certification have not been
performed by this execution. Publication and build status are recorded in
source-bound receipts, separately from mathematical completion.
